% GENERATED FILE. Edit canonical lessons, then run npm run generate:books.
% canonical-source-hash: fnv1a64:f31d3f2207a62469
% canonical-lessons: SA-C470-phupphusah, SA-C470-yakrt, SA-C470-vrkkah, SA-C470-amasayah, SA-C470-svasiti

\chapter{Inside the Body}
\label{ch:sa-inside-the-body}
\hlchaptermodality{470}

\begin{chapteropening}
\textbf{By the end of this chapter:} \emph{I can name some organs inside the body.}

Complete the last lesson of chapter 470: I can name some organs inside the body.
\end{chapteropening}

\section[\texorpdfstring{\sk{फुप्फुसः} (phupphusaḥ)}{फुप्फुसः (phupphusaḥ)}]{\sk{फुप्फुसः} (phupphusaḥ) — the lungs}

\label{lesson:SA-C470-phupphusah}



\begin{quote}
Before the new one: say the Sanskrit for a smell, then the Sanskrit for smells.
\end{quote}

\subsection*{You'll want to know: \sk{फुप्फुसः}}
\textbf{\sk{फुप्फुसः}} — \emph{phupphusaḥ} — \textquotedblleft{}the lungs\textquotedblright{}.

\begin{quote}
\sk{धूमः} \sk{फुप्फुसं} \sk{पीडयति।}
\end{quote}

\begin{quote}
\emph{dhūmaḥ phupphusaṃ pīḍayati.}
\end{quote}

\begin{quote}
Smoke harms the lungs.
\end{quote}

A masculine noun.

\begin{cousinweb}
\textbf{\sk{फुप्फुसः}} (\emph{phupphusaḥ}) MW records \emph{phupphusa} (Suśruta) and \emph{phuphusa}, \textquotedblleft{}the lungs\textquotedblright{}; MW gives no derivation, so its origin is uncertain (MW).
\end{cousinweb}

\begin{grammarlens}[title={What you've built}]
One more word for describing what you lived and felt.
\end{grammarlens}

\subsection*{Your turn}
\begin{itemize}
\raggedright
  \item \emph{Say it:} \emph{phupphusaḥ}
  \item \emph{Say it:} \emph{phupphusaḥ}, once more
  \item \emph{Say it:} \emph{dhūmaḥ phupphusaṃ pīḍayati.} — Smoke harms the lungs.
  \item \emph{Recall:} say the Sanskrit for a smell, then the Sanskrit for smells, then say \emph{phupphusaḥ} again
\end{itemize}

\begin{tcolorbox}[breakable,colback=teal!4,colframe=teal!35!black,title={Before you move on}]
What does \sk{फुप्फुसः} mean? (\textquotedblleft{}The lungs\textquotedblright{}.) Say it once more.
\end{tcolorbox}
\section[\texorpdfstring{\sk{यकृत्} (yakṛt)}{यकृत् (yakṛt)}]{\sk{यकृत्} (yakṛt) — the liver}

\label{lesson:SA-C470-yakrt}



\begin{quote}
Before the new one: say the Sanskrit for smells, then the Sanskrit for the lungs.
\end{quote}

\subsection*{You'll want to know: \sk{यकृत्}}
\textbf{\sk{यकृत्}} — \emph{yakṛt} — \textquotedblleft{}the liver\textquotedblright{}.

\begin{quote}
\sk{यकृत्} \sk{उदरे} \sk{दक्षिणतः} \sk{अस्ति।}
\end{quote}

\begin{quote}
\emph{yakṛt udare dakṣiṇataḥ asti.}
\end{quote}

\begin{quote}
The liver is on the right side of the belly.
\end{quote}

A neuter noun.

\begin{cousinweb}
\textbf{\sk{यकृत्}} (\emph{yakṛt}) MW compares Greek \emph{hēpar} and Latin \emph{jecur}, but gives no derivation within Sanskrit, so its origin there is uncertain (MW).
\end{cousinweb}

\begin{grammarlens}[title={What you've built}]
One more word for describing what you lived and felt.
\end{grammarlens}

\subsection*{Your turn}
\begin{itemize}
\raggedright
  \item \emph{Say it:} \emph{yakṛt}
  \item \emph{Say it:} \emph{yakṛt}, once more
  \item \emph{Say it:} \emph{yakṛt udare dakṣiṇataḥ asti.} — The liver is on the right side of the belly.
  \item \emph{Recall:} say the Sanskrit for smells, then the Sanskrit for the lungs, then say \emph{yakṛt} again
\end{itemize}

\begin{tcolorbox}[breakable,colback=teal!4,colframe=teal!35!black,title={Before you move on}]
What does \sk{यकृत्} mean? (\textquotedblleft{}The liver\textquotedblright{}.) Say it once more.
\end{tcolorbox}
\section[\texorpdfstring{\sk{वृक्कः} (vṛkkaḥ)}{वृक्कः (vṛkkaḥ)}]{\sk{वृक्कः} (vṛkkaḥ) — a kidney}

\label{lesson:SA-C470-vrkkah}



\begin{quote}
Before the new one: say the Sanskrit for the lungs, then the Sanskrit for the liver.
\end{quote}

\subsection*{You'll want to know: \sk{वृक्कः}}
\textbf{\sk{वृक्कः}} — \emph{vṛkkaḥ} — \textquotedblleft{}a kidney\textquotedblright{}.

\begin{quote}
\sk{मम} \sk{मित्रस्य} \sk{वृक्कः} \sk{रुग्णः} \sk{अस्ति।}
\end{quote}

\begin{quote}
\emph{mama mitrasya vṛkkaḥ rugṇaḥ asti.}
\end{quote}

\begin{quote}
My friend's kidney is diseased.
\end{quote}

A masculine noun.

\begin{cousinweb}
\textbf{\sk{वृक्कः}} (\emph{vṛkkaḥ}) MW gives \textquotedblleft{}the kidneys\textquotedblright{} in the dual, from the Atharvaveda on; MW gives no derivation, so its origin is uncertain (MW).
\end{cousinweb}

\begin{grammarlens}[title={What you've built}]
One more word for describing what you lived and felt.
\end{grammarlens}

\subsection*{Your turn}
\begin{itemize}
\raggedright
  \item \emph{Say it:} \emph{vṛkkaḥ}
  \item \emph{Say it:} \emph{vṛkkaḥ}, once more
  \item \emph{Say it:} \emph{mama mitrasya vṛkkaḥ rugṇaḥ asti.} — My friend's kidney is diseased.
  \item \emph{Recall:} say the Sanskrit for the lungs, then the Sanskrit for the liver, then say \emph{vṛkkaḥ} again
\end{itemize}

\begin{tcolorbox}[breakable,colback=teal!4,colframe=teal!35!black,title={Before you move on}]
What does \sk{वृक्कः} mean? (\textquotedblleft{}A kidney\textquotedblright{}.) Say it once more.
\end{tcolorbox}
\section[\texorpdfstring{\sk{आमाशयः} (āmāśayaḥ)}{आमाशयः (āmāśayaḥ)}]{\sk{आमाशयः} (āmāśayaḥ) — the stomach}

\label{lesson:SA-C470-amasayah}



\begin{quote}
Before the new one: say the Sanskrit for the liver, then the Sanskrit for a kidney.
\end{quote}

\subsection*{You'll want to know: \sk{आमाशयः}}
\textbf{\sk{आमाशयः}} — \emph{āmāśayaḥ} — \textquotedblleft{}the stomach\textquotedblright{}.

\begin{quote}
\sk{आमाशयः} \sk{भोजनं} \sk{पचति।}
\end{quote}

\begin{quote}
\emph{āmāśayaḥ bhojanaṃ pacati.}
\end{quote}

\begin{quote}
The stomach digests food.
\end{quote}

A masculine noun.

\begin{cousinweb}
\textbf{\sk{आमाशयः}} (\emph{āmāśayaḥ}) is built from \emph{āma}, \textquotedblleft{}raw, undigested\textquotedblright{}, and \emph{āśaya}, \textquotedblleft{}a resting place\textquotedblright{}; MW gives \textquotedblleft{}the receptacle of the undigested food, the stomach\textquotedblright{} (MW).
\end{cousinweb}

\begin{grammarlens}[title={What you've built}]
One more word for describing what you lived and felt.
\end{grammarlens}

\subsection*{Your turn}
\begin{itemize}
\raggedright
  \item \emph{Say it:} \emph{āmāśayaḥ}
  \item \emph{Say it:} \emph{āmāśayaḥ}, once more
  \item \emph{Say it:} \emph{āmāśayaḥ bhojanaṃ pacati.} — The stomach digests food.
  \item \emph{Recall:} say the Sanskrit for the liver, then the Sanskrit for a kidney, then say \emph{āmāśayaḥ} again
\end{itemize}

\begin{tcolorbox}[breakable,colback=teal!4,colframe=teal!35!black,title={Before you move on}]
What does \sk{आमाशयः} mean? (\textquotedblleft{}The stomach\textquotedblright{}.) Say it once more.
\end{tcolorbox}
\section[\texorpdfstring{\sk{श्वसिति} (śvasiti)}{श्वसिति (śvasiti)}]{\sk{श्वसिति} (śvasiti) — breathes}

\label{lesson:SA-C470-svasiti}



\begin{quote}
Before the new one: say the Sanskrit for a kidney, then the Sanskrit for the stomach.
\end{quote}

\subsection*{You'll want to know: \sk{श्वसिति}}
\textbf{\sk{श्वसिति}} — \emph{śvasiti} — \textquotedblleft{}breathes\textquotedblright{}.

\begin{quote}
\sk{शिशुः} \sk{शान्तं} \sk{श्वसिति।}
\end{quote}

\begin{quote}
\emph{śiśuḥ śāntaṃ śvasiti.}
\end{quote}

\begin{quote}
The baby breathes calmly.
\end{quote}

Past (he …): \sk{सः} \sk{शान्तं} \sk{श्वसितवान्}, \emph{saḥ śāntaṃ śvasitavān}.

\begin{cousinweb}
\textbf{\sk{श्वसिति}} (\emph{śvasiti}) comes from the root √śvas, \textquotedblleft{}to breathe\textquotedblright{}, the root of \textbf{\sk{श्वासः}} (MW).
\end{cousinweb}

\begin{grammarlens}[title={What you've built}]
One more word for describing what you lived and felt.
\end{grammarlens}

\subsection*{Your turn}
\begin{itemize}
\raggedright
  \item \emph{Say it:} \emph{śvasiti}
  \item \emph{Say it:} \emph{śvasiti}, once more
  \item \emph{Say it:} \emph{śiśuḥ śāntaṃ śvasiti.} — The baby breathes calmly.
  \item \emph{Recall:} say the Sanskrit for a kidney, then the Sanskrit for the stomach, then say \emph{śvasiti} again
\end{itemize}

\begin{tcolorbox}[breakable,colback=teal!4,colframe=teal!35!black,title={Before you move on}]
What does \sk{श्वसिति} mean? (\textquotedblleft{}Breathes\textquotedblright{}.) Say it once more.
\end{tcolorbox}
